\chapter{Run a generated input with ORCA (menu 50)}
\label{ch:menu50}
\menulabel{50}

\section{Purpose}

Hand one generated input to a local ORCA and judge the outcome.  This is
the \emph{optional execution bridge}: the core never runs calculations and
never depends on the bridge (a layered import contract keeps the
direction; dropping the subpackage removes this menu without touching
anything else), it acts only through this menu or the \code{fblockkit
submit} verb, only when explicitly asked -- and a script replay never
re-runs an engine: the run point prints one fixed ``skipped'' line
instead.

\section{What you need}

An input this toolkit generated (the bridge works on a copy, never on
your file), the full pathname of the ORCA executable (the manual requires
the complete path), a process count, the MaxCore in MB per core and a
wall-clock timeout.  The bridge creates a fresh \code{runs/<case>/}
directory per run, injects \code{\%pal}/\code{\%maxcore} only when the
input does not already carry them, launches ORCA by its full pathname
(no shell; never through \code{mpirun} itself) and monitors it -- a
heartbeat from the \code{.tmp}/\code{.gbw} timestamps, a wall-clock bound
and a stall bound, with SIGTERM-then-SIGKILL.

\section{The verdict}

The outcome is classified from three sources -- the return code, the
output text and the heartbeat -- never from the return code alone
(measured: a non-converged SCF exits 0 with an error-termination banner,
so \code{rc = 0} without the banner is \code{suspicious}, never
success):

\begin{itemize}
  \item \code{normal}: the termination banner is present;
  \item \code{orca\_error}: the engine's own error termination (the module
    name is kept);
  \item \code{input\_rejected}: the input scanner refused the file (the
    one reliable non-zero return code, 11);
  \item \code{input\_error}: a missing prerequisite file (for example the
    \code{\%moinp} gbw);
  \item \code{killed} / \code{timeout}: terminated from outside or by the
    bridge's own clocks;
  \item \code{suspicious}: the process ended without a banner.
\end{itemize}

The output and the \code{runs/<case>/run.json} ledger land beside the
copy; the ledger carries the configuration snapshot, the injection
record and the \code{.inp}/\code{.out} sha256 sums, and it is a side
product -- it never enters a byte-compared report.  Analyse the output
with menu~1 (the check-up report).

\section{Boundaries}

The bridge runs engines locally only: no queues, no ssh, no scheduler,
and no automatic re-run of a failed job (the verdict is a report, not a
repair).  Give every run a fresh directory -- measured: stale engine
residue under the same case name aborts the new run inside the engine
itself.  The generated input is the deliverable; the bridge only ever
works on a copy of it.
